\begin{table}[ht]
\centering\small
\begin{tabular}{@{}lrrrrrrrr@{}}
\toprule
kernel & $d$ & $q=1$ & $q=3$ & $q=9$ & $q=17$ & $q=33$ & $q=65$ & $v_{\min}$\\
\midrule
$\tanh$ unit gain $L=2$ & 80 & -0.04 & -1.28 & -3.26 & -5.20 & -8.86 & -22.30 & -15.3 \\
$\tanh$ unit gain $L=2$ & 160 & -0.03 & -1.26 & -3.19 & -4.95 & -7.77 & -13.18 & -23.3 \\
$\tanh$ unit gain $L=3$ & 80 & -0.05 & -1.15 & -3.06 & -5.04 & -8.78 & -22.25 & -15.3 \\
$\tanh$ unit gain $L=3$ & 160 & -0.05 & -1.14 & -2.98 & -4.78 & -7.67 & -13.13 & -23.3 \\
$\tanh$ specified gain $L=3$ & 80 & -0.07 & -1.03 & -2.67 & -4.38 & -7.82 & -20.95 & -13.9 \\
$\tanh$ specified gain $L=3$ & 160 & -0.06 & -1.02 & -2.60 & -4.15 & -6.76 & -11.89 & -21.3 \\
arcsine $L=2$ & 80 & -0.20 & -0.98 & -1.74 & -2.17 & -- & -- & -1.8 \\
arcsine $L=2$ & 160 & -0.20 & -0.98 & -1.73 & -2.15 & -2.60 & -- & -2.4 \\
\bottomrule
\end{tabular}
\caption{$\log_{10}$ of the level mass $\binom dq\lambda_q/A$ of the normalized kernel $F((1+x\cdot y)/(d+1))$,
and $\log_{10}$ of the smallest eigenvalue of the matrix $F((1+x\cdot y)/(d+1))_{x,y}$. Unit gain means
$\sigma_\ell=1$ for all $\ell$; specified gain is the family of \cite[Section~7]{nngpnote} at $L=3$.
The computation truncates the slowly decaying arcsine
coefficients at $m=600$ (omitted mass $0.021$), so arcsine entries are shown only where the omitted terms
are negligible ($q\le17$ at $d=80$, $q\le33$ at $d=160$), and the arcsine $v_{\min}$ is only a lower bound. All checks listed in Section~\ref{sec:numerics} passed for every row.}
\label{tab:spectrum}
\end{table}
